\sect{Приложение В. Среда RL и подробности реализации}{appendix-rl}
%
В настоящем приложении детально описана среда, на которой обучается и
оценивается RL-модуль предсказания темпа из~\Secref{rl}. Среда
реализована поверх симулятора, читающего размеченные исполнения из
датасета \textit{CU-Concerto-2026}, и снабжена API, совместимым с
библиотекой \texttt{gymnasium}. Один эпизод соответствует одному
шестнадцатитактному фрагменту партитуры, что даёт длину эпизода
порядка $200$ шагов агента при базовом темпе $\nu_0=120$\,bpm.

Пространство состояний $\mathcal{S}\subset\Real^{2N+2K+1}$. Координаты
включают: forward-распределение $\alpha_t\in\Real^N$, текущую
позицию $\hat\varphi(t)/N$, $K=20$ последних оценок темпа
$\tau_{t-K:t}$ (нормированных к $\nu_0$), $K$ последних значений
эмиссионной ошибки $e_{t-K:t}=-\log b_{\hat\varphi}(o_t)$ и
относительное время с начала эпизода. Для уменьшения размерности
$\alpha_t$ применяется проекция в спектральные коэффициенты
дискретного косинусного преобразования первого рода
$\alpha_t\mapsto\hat\alpha_t\in\Real^{16}$. Эмпирически такая
проекция сохраняет более $97$\,\% энергии распределения и снижает
требования к политике почти на порядок.

Пространство действий~--- одномерный отрезок $\mathcal{A}=[0{,}5;\,2{,}0]$.
Политика моделируется усечённым гауссовским распределением
$\pi_\theta(a\mid s)=\mathcal{N}_{\text{trunc}}(\mu_\theta(s),
\sigma_\theta^2)\bigr|_{[0{,}5;\,2{,}0]}$, что обеспечивает
гладкое исследование без резких выходов из физически разумного
диапазона.

Функция вознаграждения определена~\eqref{eq:rl-reward} и состоит из
трёх компонентов. Первый~--- штраф за рассогласование
$|t^{\text{render}}_t-t^{\text{perf}}_t|$~--- измеряется в
миллисекундах, нормирован на $50$\,мс (поскольку именно это значение
является порогом восприятия рассинхронизации в музыкальном
ансамбле, см.~\Secref{introduction}). Второй компонент~---
$\mathcal{L}_{\text{align}}(t)$~--- равен расстоянию по партитуре
между ground-truth и текущей оценкой $\hat\varphi$, нормированному
на длину фрагмента. Третий компонент~--- штраф гладкости
$(a_t-a_{t-1})^2$~--- ограничивает резкие переключения темпа,
которые акустически воспринимаются как «дёрганье» оркестра.

Симулятор работает в двух режимах. В \textit{train-режиме} к
номинальному темпу записанного исполнения применяются случайные
возмущения, эмулирующие индивидуальное rubato: умножение
длительностей на гауссовский шум $\mathcal{N}(1,\sigma_\tau^2)$,
случайные локальные ритардандо и ачелерандо в окне восьми тактов и
вставка пропусков с вероятностью $0{,}5$\,\% на ноту. В
\textit{eval-режиме} симулятор воспроизводит ровно записанное
исполнение без возмущений; этот режим используется для
итоговой оценки на отложенной выборке.

\par\addvspace{2pt}
{\centering
\refstepcounter{table}\tablab{rl-env}%
\footnotesize
\begin{tabular}{@{}lc@{}}
  \toprule
  Параметр среды                           & Значение \\
  \midrule
  Длина эпизода                            & $16$ тактов ($\sim 200$ шагов) \\
  Шаг агента                               & $10$\,мс \\
  Размерность состояния                    & $73$ ($16+20+20+1+16$) \\
  Размерность действия                     & $1$ \\
  Возмущение темпа в train-режиме          & $\sigma_\tau\in[0{,}02;\,0{,}12]$ \\
  Вероятность пропуска ноты                & $5\!\times\!10^{-3}$ \\
  Вероятность повтора ноты                 & $2\!\times\!10^{-3}$ \\
  Норма для штрафа рассогласования         & $50$\,мс \\
  Норма для штрафа выравнивания            & длина фрагмента \\
  \bottomrule
\end{tabular}\par
\vspace{2pt}
\tablename~\thetable. Параметры симулятора, реализующего среду
для обучения и оценки RL-модуля предсказания темпа.\par}
\addvspace{4pt}

Реализация написана на Python~3.11 с использованием библиотеки
\texttt{stable-baselines3} в части PPO и \texttt{torch} версии $2{,}3$
в части политики. Все вычисления, не связанные с обучением, проводятся
на CPU; обучение PPO~--- на одной GPU NVIDIA~RTX~3060. Полная сходимость
достигается за $\sim 12$ часов при $5\!\times\!10^{6}$ шагах среды;
behavioural cloning требует около $25$ минут на тех же исходных
данных. Соответствующие веса политики сохраняются в формате
\texttt{state\_dict} и загружаются в продакшн-конвейер
(\Secref{realtime}) как обычный PyTorch-чекпойнт.
